\section{\texorpdfstring{More tactics: \href{https://lean-lang.org/doc/reference/latest/Tactic-Proofs/Tactic-Reference/}{\texttt{constructor}}, \href{https://lean-lang.org/doc/reference/latest/Tactic-Proofs/Tactic-Reference/}{\texttt{cases}}, \href{https://lean-lang.org/doc/reference/latest/Tactic-Proofs/Tactic-Reference/}{\texttt{induction}}, \href{https://lean-lang.org/doc/reference/latest/Tactic-Proofs/Tactic-Reference/}{\texttt{unfold}}, \href{https://lean-lang.org/doc/reference/latest/Tactic-Proofs/Tactic-Reference/}{\texttt{simp}}}{More tactics: constructor, cases, induction, unfold, simp}}\label{sec:05-tactics:04-more-tactics:1}

None of \texttt{intro}/\texttt{exact}/\texttt{apply}/\texttt{rw} from the previous section says anything about a goal with two independent parts to prove at once, like \(P \wedge Q\): proving it needs a proof of \(P\) \emph{and} a proof of \(Q\), two separate obligations from one goal.

\begin{lstlisting}
theorem and_example (P Q : Prop) (hp : P) (hq : Q) : P (*$\wedge$*) Q := by
  constructor
  (*$\cdot$*) exact hp
  (*$\cdot$*) exact hq
\end{lstlisting}

\href{https://lean-lang.org/doc/reference/latest/Tactic-Proofs/Tactic-Reference/}{\texttt{constructor}} is exactly the tactic for this: it invokes the introduction rule of the goal's type, splitting \(P \wedge Q\) into the two goals \(P\) and \(Q\) (the \texttt{·} focus dot addresses each one in turn), and more generally, for any \texttt{structure}, into one goal per field. This reflects the universal property of the product, ``a map into \(P \times Q\) is a pair of maps,'' so proving \(P\) and proving \(Q\) is enough.

\begin{quote}
Read more. \hyperref[sec:04-propositions-and-proofs:02-logic-recap:1]{Chapter 4, Section 2} states the introduction/elimination rules for every connective by name, if ``introduction rule''/``elimination rule'' as general terms are new.
\end{quote}

\texttt{constructor} builds a conjunction; the dual problem is \emph{consuming} a disjunction already in hand. A hypothesis \(h : P \vee Q\) gives no term directly usable by \texttt{exact}, since which of \(P\), \(Q\) actually holds is not known, only that one of them does.

\begin{lstlisting}
theorem or_comm_ex {P Q : Prop} (h : P (*$\vee$*) Q) : Q (*$\vee$*) P := by
  cases h with
  | inl hp => exact Or.inr hp
  | inr hq => exact Or.inl hq
\end{lstlisting}

\href{https://lean-lang.org/doc/reference/latest/Tactic-Proofs/Tactic-Reference/}{\texttt{cases}} splits the proof into one branch per way \(h\) could have been built, \texttt{inl hp : P} or \texttt{inr hq : Q}, and each branch supplies its own proof of the goal. This is the \emph{elimination} rule of a coproduct (the reading of \(\vee\) as a coproduct in \hyperref[sec:04-propositions-and-proofs:05-and-or-not:1]{Chapter 4, Section 5}): to prove anything from \(P \sqcup Q\), it suffices to prove it from each summand, the case analysis \(\iota_1\)/\(\iota_2\). \texttt{cases} works the same way on any inductive value, not only a disjunction hypothesis, splitting into one goal per constructor.

\texttt{cases} handles a value built by \emph{one} of finitely many constructors, but \texttt{Nat} has a constructor, \texttt{succ}, that refers back to a smaller \texttt{Nat}. Splitting on \texttt{succ k} alone leaves a goal about \texttt{k} with nothing yet known about it, unlike the finitely many closed cases \texttt{cases} produces elsewhere.

\begin{lstlisting}
theorem add_zero_left (n : Nat) : 0 + n = n := by
  induction n with
  | zero => rfl
  | succ k ih =>
    show 0 + (k + 1) = k + 1
    rw [Nat.add_succ, ih]
\end{lstlisting}

\href{https://lean-lang.org/doc/reference/latest/Tactic-Proofs/Tactic-Reference/}{\texttt{induction}} is \texttt{cases} strengthened for exactly this: the \texttt{succ k} branch additionally supplies \texttt{ih}, a proof of the goal \emph{for \texttt{k}}, the mathematical principle of induction made available as a hypothesis.

\[
P(0) \;\land\; \big(\forall k,\ P(k) \to P(k+1)\big) \;\implies\; \forall n,\ P(n)
\]

\texttt{ih} is exactly \(P(k)\), available to prove \(P(k+1)\).

None of the tactics so far can see through a name like \texttt{isZero} to what it actually means; \texttt{rw} needs an equation to rewrite with, and \texttt{isZero 0} is not stated as one.

\begin{lstlisting}
def isZero (n : Nat) : Prop := n = 0

theorem isZero_zero : isZero 0 := by
  unfold isZero
  -- Goal after unfold: 0 = 0
  rfl
\end{lstlisting}

\href{https://lean-lang.org/doc/reference/latest/Tactic-Proofs/Tactic-Reference/}{\texttt{unfold}} replaces the defined name with its definiens, exposing \(\mathrm{isZero}(0) = (0 = 0)\). This uses a \emph{definitional} equality: \(\mathrm{isZero} := (n \mapsto n = 0)\) means the two are interchangeable by definition (like unwinding ``\(n\) is even'' to ``\(\exists k,\ n = 2k\)'' in a proof). So the goal becomes the tautology \(0 = 0\), closed by reflexivity.

\subsection{\texorpdfstring{An aside on transparency: \texttt{def} vs.~\texttt{abbrev} vs.~\texttt{opaque}}{An aside on transparency: def vs.~abbrev vs.~opaque}}\label{sec:05-tactics:04-more-tactics:2}

\texttt{unfold isZero} above had to name \texttt{isZero} explicitly, the same way \texttt{rw} names an equation explicitly, because that is exactly what \texttt{unfold} is for. This is a good moment to name the promise from \hyperref[sec:01-basics:02-def-let-implicit:1]{Chapter 1, Section 2}. \texttt{def} is not the only way to introduce a definition, and the alternatives differ precisely in how much they need to be \emph{told about} like this, versus being seen through automatically.

\begin{itemize}
\tightlist
\item
  \textbf{\texttt{def} (semi-reducible, the default).} Requires an explicit instruction to unfold, \texttt{unfold} by name in a tactic proof, exactly as above. This matches the traceability this book has been deliberately practicing throughout this chapter, nothing unfolds silently, and every step names what justified it.
\item
  \textbf{\texttt{abbrev} (reducible, automatically \texttt{@[reducible, inline]}).} Declares the definition \emph{notational}, not merely equal, so the elaborator of Lean and automated machinery treat every occurrence as if the body had been written out directly, with no \texttt{unfold} step to ask for. The clean demonstration of this has to wait for tools that actually \emph{search} through definitions automatically. Typeclass resolution (\hyperref[sec:06-rigor-check:01-structure-vs-class:1]{Chapter 6}) searches at reducible transparency by default, so an \texttt{abbrev} (or a \texttt{def} tagged \texttt{@[reducible]}, as \texttt{isPrime} was in \hyperref[sec:04-propositions-and-proofs:06-quantifiers:1]{Chapter 4}) is visible to that search without being unfolded by hand first, while a plain \texttt{def} generally is not.
\item
  \textbf{\texttt{opaque}.} The opposite extreme, not unfoldable at all, by \texttt{unfold} or anything else, even though it still has a definition somewhere. Useful for genuinely hiding an implementation and exposing only the properties proved about it, the Lean equivalent of citing a chosen but unspecified witness (``let \(c\) be \emph{some} element of the nonempty set \(S\)'') and reasoning only from what is known about it, never from how it was built.
\end{itemize}

\begin{mathreading}

This maps onto a distinction already made informally, if rarely by name, in ordinary mathematical writing. A \texttt{def} is a definition cited by name in a proof, ``by definition of \(f\), \ldots{}'', an explicit, visible step. An \texttt{abbrev} is pure notation, the kind silently unfolded without comment, writing \(2n\) in place of ever having introduced a name \(f(n) := 2n\) at all. An \texttt{opaque} value is an axiomatized or black-boxed object, reasoned about only through its stated properties, never its construction. The three transparency levels of Lean are exactly these three mathematical habits, made explicit and machine-checked instead of left to convention.

\end{mathreading}

Every proof in this chapter so far has named, by hand, exactly which lemma or equation justifies each step: \texttt{Nat.add\_\allowbreak{}succ}, \texttt{ih}, an \texttt{unfold}d definition. Once the same handful of rewrites recur across many similar goals, naming each one individually stops teaching anything new and becomes pure repetition.

\begin{lstlisting}
theorem simp_example (n : Nat) : n + 0 = n := by
  simp
\end{lstlisting}

\href{https://lean-lang.org/doc/reference/latest/Tactic-Proofs/Tactic-Reference/}{\texttt{simp}} automates exactly this: it searches a library of known ``simplification'' lemmas and applies as many as fire, in one step. The cost is exactly the traceability just built up by hand: \texttt{simp} hides \emph{which} facts were used and \emph{why} the proof works, which is bad for learning something for the first time. \textbf{This book therefore avoids \texttt{simp} and \texttt{rfl}-as-a-shortcut wherever the point is to understand the proof}, using explicit \texttt{rw} and \texttt{induction} with a fully spelled-out base case and inductive step instead. \texttt{simp} is a tool for \emph{later}, once what it would have done by hand is already understood, not a substitute for doing it by hand the first time.

\begin{quote}
Read more. \hyperref[sec:13-working-efficiently:03-simp:1]{Chapter 13, ``\texttt{simp}, now that you understand what it replaces''} covers when it is the \emph{right} efficient choice, once every step no longer needs to be spelled out.
\end{quote}

\hyperref[sec:05-tactics:03-reading-failures:1]{← Reading failures} \textbar{} \hyperref[sec:05-tactics:00-index:1]{Index} \textbar{} \hyperref[sec:05-tactics:05-worked-example:1]{Next: Worked example →}
